\documentclass[11pt,a4paper]{article}

\usepackage[utf8]{inputenc}
\usepackage[T1]{fontenc}
\usepackage{amsmath,amssymb,amsthm,mathtools}
\usepackage{bm}
\usepackage{graphicx}
\usepackage[margin=1.1in]{geometry}
\usepackage{enumitem}
\usepackage{booktabs}
\usepackage{float}
\usepackage{natbib}
\usepackage{microtype}
\usepackage{lmodern}
\usepackage[version=4]{mhchem}
\usepackage[colorlinks=true,linkcolor=blue,citecolor=blue,urlcolor=blue]{hyperref}
\usepackage{cleveref}
\crefname{postulate}{Postulate}{Postulates}
\Crefname{postulate}{Postulate}{Postulates}

\theoremstyle{plain}
\newtheorem{theorem}{Theorem}[section]
\newtheorem{lemma}[theorem]{Lemma}
\newtheorem{proposition}[theorem]{Proposition}
\newtheorem{corollary}[theorem]{Corollary}

\theoremstyle{definition}
\newtheorem{definition}[theorem]{Definition}
\newtheorem{postulate}{Postulate}
\newtheorem{example}[theorem]{Example}

\theoremstyle{remark}
\newtheorem{remark}[theorem]{Remark}

% --- notation -----------------------------------------------------------
\newcommand{\kB}{k_{\mathrm{B}}}
\newcommand{\Univ}{\mathcal{U}}
\newcommand{\Item}{\mathcal{I}}
\newcommand{\thick}{\beta}
\newcommand{\floor}{\thick_{\min}}
\newcommand{\Cat}{\mathfrak{C}}
\newcommand{\Compl}{\gamma}
\newcommand{\Scat}{S_{\mathrm{cat}}}
\newcommand{\Skin}{S_{\mathrm{kin}}}
\newcommand{\dC}{d_{\mathrm{C}}}
\newcommand{\Depth}{\mathcal{M}}
\newcommand{\taup}{\tau_{\mathrm{p}}}
\newcommand{\nufloor}{\nu_{\mathrm{floor}}}
\newcommand{\kcat}{k_{\mathrm{cat}}}
\newcommand{\KM}{K_{\mathrm{M}}}
\newcommand{\Keq}{K_{\mathrm{eq}}}
\newcommand{\Enz}{E}
\newcommand{\Sub}{S}
\newcommand{\prov}{\mathcal{P}}
\newcommand{\rel}{\varrho}

\title{\textbf{Why Enzymes Exist:\\
Catalysis as Categorical Provision Upon Contact}}

\author{Kundai Farai Sachikonye\\
\small Technical University of Munich, School of Life Sciences\\
\small \texttt{kundai.sachikonye@wzw.tum.de}}

\date{\today}

\begin{document}
\maketitle

\begin{abstract}
\noindent
Catalysis is conventionally described but not derived. Transition-state theory
tells us that a catalyst lowers an activation barrier; it does not tell us why
such objects exist, why they are reusable, why they are specific, or why
specificity is bounded above as well as below. We derive catalysis from two
primitives. The first is that individuation is not free: distinguishing an item
from its complement takes finite time, during which the complement evolves, so
every boundary carries irreducible thickness $\thick>0$. The second is that a
\emph{category} is the unit of a step---the structure that permits an invariant
to persist across a transition in which nothing numerically survives. We prove
these primitives are independent: the first drives systems toward configurations
of minimum boundary cost, but it cannot account for reuse, because the
uncertainty a process resolves is consumed in resolving it. Reuse requires
something that is identically the same across occasions while each occasion is a
distinct event, and that is what a category is and residue is not. From this we
obtain the definition: \emph{a catalyst is an item that, upon contact, increases
the number of completable categories available to a transition, and is not
consumed because categorical provision requires no free energy.} Nine results
follow as corollaries rather than assumptions: catalysts are not consumed;
equilibrium constants are unaltered; the Haldane relation is recovered; rate
enhancement is exponential in the number of apertures traversed with a single
framework-wide constant; specificity is bounded on both sides by the joint
requirement of provision and release, giving a \emph{specificity window};
competitive inhibition is contact without provision; mechanism-based
inactivation is provision without release; allosteric regulation is
modification of the provider's contact chain; and catalysis is not a Maxwell
demon, because categories are configurational and the demon's criterion is
kinetic, quantities which depend on disjoint variables. Because the account
never requires that the provider be a protein, it resolves the origin-of-life
regress in which enzymes require proteins and proteins require nucleic acid
templates: the first catalysts were structures whose delocalised electronic
motifs supplied categories on contact and were unchanged by their use. We give
an abstract system exhibiting the same mechanics with no chemistry in it,
establishing that the categorical primitive is not a chemical postulate. All
claims are stated as falsifiable predictions and are accompanied by
computational validation against public reaction and kinetics databases.
\end{abstract}

\tableofcontents

%==============================================================================
\part{The Problem}
%==============================================================================

\section{Introduction}
\label{sec:introduction}

\subsection{What is not explained by the standard account}
\label{sec:not-explained}

The standard account of enzyme catalysis is transition-state stabilisation
\citep{pauling1946,eyring1935}. An enzyme binds the transition state more
tightly than the ground state; the activation free energy falls; the rate rises
\citep{schramm2011,warshel2006}. This account is correct as far as it goes, and
it is quantitatively successful. It is also, as an explanation, incomplete in
five specific ways, and the incompleteness is not a matter of missing detail but
of missing derivation.

\paragraph{First: why does anything play this role at all?}
Transition-state stabilisation describes what an enzyme does once we are given
an enzyme. It does not explain why objects of this kind should exist, or why
matter organised into cells should invariably produce them. To say that enzymes
exist because they are useful is to give a selective rather than a physical
answer, and it presupposes the existence of the entities selection acts upon.

\paragraph{Second: why is the catalyst reusable?}
This is the question we shall find to be decisive. A catalyst participates in a
reaction and emerges unaltered, ready to participate again. On any account in
which catalysis consists in the resolution of some uncertainty, or the traversal
of some barrier, or the payment of some cost, the first turnover ought to
exhaust the thing that made it possible. That it does not is normally recorded
as a definitional stipulation---``a catalyst is not consumed''---rather than
derived. We shall show that the requirement of reuse is strong enough to force a
second primitive into the theory.

\paragraph{Third: why is a catalyst specific?}
Specificity is usually explained by complementarity \citep{fischer1894,
koshland1958}: the substrate fits the site. But complementarity explains only
the lower bound---why a badly matched substrate is not turned over. It does not
explain the upper bound. Nothing in the lock-and-key picture forbids a lock so
exquisitely matched to its key that the two never separate. Yet such objects are
not enzymes; they are, at best, single-use reagents. Real enzymes occupy a
bounded region of specificity, and the bound above is not explained.

\paragraph{Fourth: why do inhibitors fall into the classes they do?}
Competitive inhibitors occupy the site and do nothing. Mechanism-based
inactivators are turned over once and destroy the enzyme in the process
\citep{abeles1976,silverman1995}. These are treated as separate pharmacological
phenomena. On the account developed here they are the two ways the catalytic
requirement can fail, and they are therefore predictions rather than
observations.

\paragraph{Fifth: what were the first catalysts?}
Life requires enzymes. Enzymes, in extant biology, are proteins. Proteins
require templated synthesis, hence nucleic acids, hence a replication apparatus
that is itself enzyme-dependent \citep{eigen1971,gilbert1986,orgel2004}. The
regress is well known and the standard resolutions---ribozymes
\citep{kruger1982,guerrier1983}, surface metabolism \citep{wachtershauser1988},
autocatalytic sets \citep{kauffman1986}---all proceed by exhibiting a non-protein
catalyst. What is missing is a criterion that tells us what makes something a
catalyst independently of what it is made of. Without such a criterion the
resolutions are existence proofs without a theory.

\subsection{What this paper does}
\label{sec:what-this-does}

We derive catalysis from two primitives, prove they are independent, and obtain
the five missing explanations as consequences.

The first primitive (\cref{sec:individuation}) is that individuation costs
something irreducibly. To be a distinct item is to be distinguished from
everything one is not; establishing that distinction takes non-zero time; during
that time the complement changes; so the boundary is not a clean cut but a layer
of finite thickness $\thick > 0$. We derive this rather than assume it, and we
derive it in a setting where the bound is uniform, which is the technically
delicate point.

The second primitive (\cref{sec:category}) is the \emph{category}. We motivate it
by an argument from reuse: what a catalyst must have, in order to be usable
twice, is an invariant that is numerically the same on both occasions while the
occasions are genuinely distinct events. We show that the first primitive cannot
supply this, and that the structure which can is what we shall call a category.
We then give the measure---categorical entropy---and prove the property that
makes catalysis possible: completing a category increases entropy without
consuming free energy.

\Cref{sec:abstract-system} exhibits an abstract system, containing no chemistry
whatever, which has both primitives and in which the second is demonstrably not
reducible to the first. This is included because the central claim of the paper
is a claim about what kind of thing a catalyst is, and such a claim should not
depend on the accident that our examples are molecules.

\Cref{part:catalysis} contains the definition and its consequences.
\Cref{part:predictions} states the falsifiable predictions and the validation
protocol. \Cref{part:origins} applies the definition to the origin-of-life
regress.

\subsection{Conventions}
\label{sec:conventions}

Throughout, $\Univ$ denotes the totality under discussion and $\Item$ an item
individuated within it. We write $\thick(A,B)$ for the boundary thickness
between $A$ and $B$ and $\floor$ for its lower bound. $\kB$ is Boltzmann's
constant. Logarithms are natural unless subscripted. We use $\Omega$ for a count
of microstates in the Boltzmann sense \citep{boltzmann1877,gibbs1902} and
$|\Compl(t)|$ for a count of completed categories. Standard enzyme-kinetic
notation follows \citet{cornish1979}: $\kcat$, $\KM$, $\Keq$, with
$\kcat/\KM$ the specificity constant.

We are careful throughout to distinguish results that are unconditional from
those that hold under stated hypotheses, and we mark the latter explicitly.

%==============================================================================
\part{Foundations}
%==============================================================================

\section{Individuation Is Not Free}
\label{sec:individuation}

\subsection{Existence by distinction}

Physics customarily begins with objects already individuated: particles,
systems, states. The process by which something comes to be a distinct
something is not usually part of the formalism. We make it part of the
formalism, because catalysis turns out to be a fact about that process.

\begin{postulate}[Existence by distinction]
\label{post:negation}
An item $\Item$ exists as a distinct entity if and only if there is a
well-defined complement $\Univ \setminus \Item$ with
$\Item \cap (\Univ\setminus\Item) = \varnothing$ and
$\Item \cup (\Univ\setminus\Item) = \Univ$. The existence of $\Item$ consists in
the distinction $\Item \neq \Univ\setminus\Item$, and not in any property
$\Item$ possesses independently of it.
\end{postulate}

This is not a metaphysical flourish. It is the observation that every physical
predicate we assign to an item---its mass, its charge, its position---is
assigned by contrast with what does not have it. \Cref{post:negation} says only
that we should take that contrast as constitutive rather than incidental.

\begin{postulate}[Distinction takes time]
\label{post:duration}
Establishing the distinction $\Item \neq \Univ\setminus\Item$ requires a finite,
non-zero interval $\taup > 0$. During $\taup$ the complement
$\Univ\setminus\Item$ evolves.
\end{postulate}

\begin{remark}
\Cref{post:duration} is forced rather than chosen. Suppose $\taup = 0$. Then the
state of $\Univ$ immediately before the distinction and immediately after are
the same state. But if nothing about $\Univ$ differs, no distinction has been
drawn---the supposed individuation made no difference, and an individuation that
makes no difference is not one. Hence $\taup > 0$. The second clause is then
automatic for any $\Univ$ that is not static.
\end{remark}

\subsection{The boundary has thickness}

\begin{definition}[Boundary thickness]
\label{def:thickness}
Let $A$ be individuated within $\Univ$ over the interval $[t_0, t_0+\taup]$. The
\emph{boundary thickness} $\thick(A)$ is the measure of the set of states which,
at the completion of the individuation, belong determinately neither to $A$ nor
to $\Univ\setminus A$.
\end{definition}

\begin{theorem}[Positivity of thickness]
\label{thm:floor}
For any individuation carried out by a process satisfying
\cref{post:negation,post:duration}, $\thick(A) > 0$.
\end{theorem}

\begin{proof}
Let the individuation begin at $t_0$ and complete at $t_0+\taup$ with
$\taup>0$ by \cref{post:duration}. Consider a state $s$ that lies in
$\Univ\setminus A$ at $t_0$ and in the region assigned to $A$ at $t_0+\taup$, or
conversely. Such states exist whenever $\Univ$ is non-static on the timescale
$\taup$, which \cref{post:duration} guarantees. For such $s$ the assignment made
at the start of the operation and the assignment appropriate at its completion
disagree, so $s$ is not determinately in $A$ nor determinately in
$\Univ\setminus A$. The set of such states is non-empty, and $\thick(A)$ is its
measure, hence positive.
\end{proof}

\Cref{thm:floor} gives positivity for each individuation separately. For the
results of \cref{part:catalysis} we require something stronger: a bound that is
uniform over the individuations a given system can perform. Positivity of each
member of a family does not imply positivity of the infimum---the numbers
$1/n$ are all positive and have infimum zero---so the uniform statement requires
a hypothesis that bounds refinement.

\begin{definition}[Resolution-bounded system]
\label{def:res-bounded}
A system is \emph{resolution-bounded} if the individuations it can perform are
drawn from a finite set, or more generally if there is a finest individuation it
can carry out, of thickness $\mu_{\min} > 0$.
\end{definition}

\begin{theorem}[Uniform floor]
\label{thm:uniform-floor}
In a resolution-bounded system there exists $\floor > 0$ such that
$\thick(A) \geq \floor$ for every individuated $A$. One may take
$\floor = \mu_{\min}$.
\end{theorem}

\begin{proof}
By \cref{def:res-bounded} the set of attainable thicknesses has a least element
$\mu_{\min}$, or is bounded below by one. By \cref{thm:floor} every attainable
thickness is positive, so $\mu_{\min}>0$. The minimum of a finite set of
positive numbers is positive; where the set is infinite the hypothesis of
\cref{def:res-bounded} supplies the bound directly. Hence
$\thick(A) \geq \mu_{\min} = \floor > 0$ for all individuated $A$.
\end{proof}

\begin{remark}[Where the hypothesis is needed and where it is not]
\label{rem:hypothesis}
\Cref{thm:floor} is unconditional. \Cref{thm:uniform-floor} is conditional on
\cref{def:res-bounded}, and the condition is not cosmetic: a system able to
refine without limit can have thicknesses approaching zero, and for such a
system no uniform floor exists. Every physical system we shall consider is
resolution-bounded, because the number of distinguishable cells in a bounded
region of phase space is finite. We flag the dependence because results
downstream inherit it.
\end{remark}

\subsection{Contact}

\begin{definition}[Contact]
\label{def:contact}
Items $A$ and $B$ are in \emph{contact} when their mutual boundary thickness
attains the floor:
\begin{equation}
\thick(A,B) = \floor .
\label{eq:contact}
\end{equation}
\end{definition}

Contact is the closest two items can come while remaining distinct. Below
$\floor$ they are not two items; at $\floor$ they are two items maximally
mutually defined.

\begin{proposition}[Local distinction is cheaper than global]
\label{prop:local-cheaper}
Maintaining $A$ as distinct from a specific finite partner $B$ requires boundary
content bounded by $\thick(A,B)$; maintaining $A$ as distinct from the whole of
$\Univ$ requires boundary content bounded below by $\thick(A,\Univ)$, and
$\thick(A,B) \leq \thick(A,\Univ)$ whenever $B \subseteq \Univ\setminus A$.
\end{proposition}

\begin{proof}
The distinction of $A$ from $\Univ\setminus A$ entails its distinction from
every $B \subseteq \Univ\setminus A$, so the set of indeterminate states at the
$A$--$B$ boundary is a subset of that at the $A$--$\Univ$ boundary. Measures
respect inclusion.
\end{proof}

\begin{corollary}[Contact is favoured]
\label{cor:contact-favoured}
Absent constraints preventing it, items evolve toward contact, because contact
minimises the boundary content required to maintain mutual distinctness.
\end{corollary}

\Cref{cor:contact-favoured} is a tendency, not a force. Nothing pulls the items
together. Configurations of lower total boundary cost are simply the ones that
persist, in the same sense in which a system samples its accessible states.

\begin{definition}[Contact chain]
\label{def:chain}
A \emph{contact chain} is a sequence $A_1 - A_2 - \cdots - A_n$ in which each
consecutive pair is in contact.
\end{definition}

\begin{proposition}[Chains sharpen without global cost]
\label{prop:chain}
Let $A_1-\cdots-A_n$ be a contact chain and let $A_0$ come into contact with
$A_1$. Then each of $A_1,\dots,A_n$ is thereby distinguished more finely---each
acquires an additional specific partner against which it is defined---while no
$\thick(A_i, \Univ)$ increases.
\end{proposition}

\begin{proof}
The new contact establishes $\thick(A_0,A_1) = \floor$ without altering any
existing $\thick(A_i,A_{i+1})$, each of which is already at $\floor$ by
\cref{def:chain}. The context of each $A_i$ now includes one further specific
partner, which refines the distinctions available to it. No global boundary is
modified, since $A_0$ was already individuated within $\Univ$.
\end{proof}

\subsection{What the first primitive gives, and what it does not}
\label{sec:residue-limits}

\Cref{post:duration} has a consequence beyond thickness. Because individuation
takes time and the complement moves, each individuation leaves an
\emph{unresolved remainder}: content that was one thing when the operation began
and another when it completed, and which the operation therefore failed to
assign. Call this the residue.

\begin{proposition}[Residue drives boundary minimisation]
\label{prop:residue-drives}
Accumulated residue is monotone: no operation removes it, since removal would
require re-assigning content whose assignment was indeterminate, which is itself
an individuation generating fresh residue. Systems therefore evolve
irreversibly, and \cref{cor:contact-favoured} gives the direction: toward
configurations of lower total boundary cost.
\end{proposition}

This is a genuine dynamical principle, and it is all the first primitive gives.
It is not enough, and the way in which it fails is the hinge of this paper.

\begin{proposition}[Residue cannot account for reuse]
\label{prop:no-reuse}
Let $X$ be an item whose contribution to a transition consists solely in the
residue that transition generates. Then $X$ cannot support a second transition
of the same kind.
\end{proposition}

\begin{proof}
Residue is generated by the resolution of an indeterminacy. Once the transition
has occurred, the indeterminacy it resolved is resolved: by
\cref{prop:residue-drives} it is not restored, since restoration is itself an
individuation and generates its own residue rather than undoing the prior one.
The condition under which $X$ made the first transition possible therefore does
not obtain at the second occasion. Hence $X$ is, in the relevant respect,
consumed.
\end{proof}

\Cref{prop:no-reuse} is the crux. A catalyst is by observation reusable. If
everything a catalyst supplies were residue, catalysts would be single-use, and
the phenomenon we are trying to explain would not exist. Something else must be
supplied---something which is numerically the same on the second occasion as on
the first, while the two occasions are distinct events. We now say what that is.

\section{Category: The Unit of a Step}
\label{sec:category}

\subsection{The requirement}

Consider what \cref{prop:no-reuse} demands. We require a structure $\Cat$ with
two properties that are in apparent tension:

\begin{enumerate}[label=(\roman*),leftmargin=2.5em]
\item \textbf{Invariance.} $\Cat$ is the same at the second use as at the first.
Were it different, what made the first turnover possible would not be what makes
the second possible, and the reuse would be a coincidence rather than a
mechanism.
\item \textbf{Difference.} The first use and the second use are distinct events.
Were they not distinct, no second turnover has occurred.
\end{enumerate}

An item cannot be numerically identical across a change in which it participates
and also be what changed. But it can be the \emph{same category} while the
\emph{completions} of that category are distinct. This is the ordinary structure
of a type and its tokens, imported into physics as a structural rather than a
linguistic fact.

\begin{definition}[Category]
\label{def:category}
A \emph{category} is a distinction that can be completed. It is specified
independently of any particular completion; a \emph{completion} is an occasion
on which the distinction is drawn. The category is invariant across its
completions; the completions are distinct events.
\end{definition}

\begin{remark}[Why this is a second primitive]
\label{rem:second-primitive}
A category is not constructed from residue and cannot be. Residue is what a
completion leaves behind and is consumed in the leaving; a category is what
makes completions of a given kind possible and survives them. The two are
related as a capacity is related to its exercises. \Cref{prop:no-reuse} shows
the capacity is not recoverable from the exercises. We therefore take
\cref{def:category} as primitive alongside
\cref{post:negation,post:duration}, and \cref{sec:abstract-system} exhibits a
system in which the independence can be checked directly.
\end{remark}

\begin{remark}[Countability]
\label{rem:countability}
Completions are countable: at any time a finite number have occurred, and
$|\Compl(t)|$ denotes that number. The categories available for completion are
not assumed countable, and nothing below requires that they be. All measures are
defined on completions.
\end{remark}

\subsection{Categorical entropy}

\begin{definition}[Categorical entropy]
\label{def:scat}
For a system whose completed categories at time $t$ number $|\Compl(t)|$, the
\emph{categorical entropy} is
\begin{equation}
\Scat(t) = \kB \ln |\Compl(t)| .
\label{eq:scat}
\end{equation}
\end{definition}

\begin{definition}[Kinetic entropy]
\label{def:skin}
The \emph{kinetic entropy} is the familiar quantity measuring the distribution
of energy over accessible microstates,
\begin{equation}
\Skin = -\kB \sum_i p_i \ln p_i .
\label{eq:skin}
\end{equation}
\end{definition}

\begin{theorem}[Completion increases entropy]
\label{thm:completion-increases}
Let $\Omega(C)$ be the number of microstates compatible with $C$ completed
categories. Then $\Omega$ is strictly increasing in $C$, and consequently
\begin{equation}
\frac{d\Scat}{dC} = \kB \frac{1}{\Omega(C)}\frac{d\Omega}{dC} > 0 .
\label{eq:dscat}
\end{equation}
\end{theorem}

\begin{proof}
A completed category is a distinction the system did not previously draw: a new
respect in which its configurations differ from one another. Configurations
indistinguishable before the completion are distinguishable after it. Hence
every microstate accessible at $C$ remains accessible at $C+\Delta C$, and
additional microstates---those differing only in the newly drawn
distinction---become accessible. So $\Omega(C+\Delta C) > \Omega(C)$ for
$\Delta C > 0$. Applying $S = \kB\ln\Omega$ and differentiating gives
\cref{eq:dscat}.
\end{proof}

\subsection{The property that makes catalysis possible}

\begin{theorem}[Completion at zero free energy]
\label{thm:zero-free-energy}
Categorical completion does not require free energy. There exist transitions
with $\Delta F = 0$, $Q = 0$ and $W = 0$ which nevertheless complete a category
and increase $\Scat$ by at least $\kB \ln 2$.
\end{theorem}

\begin{proof}
Consider a system at thermal equilibrium at temperature $T$, so that
$\Delta F = 0$ for every spontaneous process. Let its internal configuration be
specified by a vector of mode occupations $\mathbf{v} = (n_1,\dots,n_M)$, and
consider a transition $\mathbf{v} \to \mathbf{v}'$ that redistributes occupation
among modes at fixed total. By equipartition the mean energy per mode is fixed
by $T$, so $\langle E(\mathbf{v})\rangle = \langle E(\mathbf{v}')\rangle$: no
net energy flows and $Q=0$. No macroscopic displacement occurs, so $W=0$. Hence
\begin{equation}
\Delta F = \Delta U - T\Delta S = 0 .
\end{equation}
But $\mathbf{v}$ and $\mathbf{v}'$ are distinguishable configurations, so the
transition completes a category not previously completed. By
\cref{thm:completion-increases} this increases $\Scat$, and the minimal increase
associated with drawing one binary distinction is $\kB\ln 2$.
\end{proof}

\begin{corollary}[Provision is free]
\label{cor:provision-free}
An item which supplies completable categories to another item, without itself
undergoing a transition that consumes free energy, incurs no free-energy cost in
doing so, and may repeat the provision indefinitely.
\end{corollary}

\begin{proof}
Immediate from \cref{thm:zero-free-energy}: the completions the item enables are
$\Delta F = 0$ transitions, and by \cref{def:category} the category itself is
invariant across its completions. Nothing about the provider is spent.
\end{proof}

\Cref{cor:provision-free} is the resolution of \cref{prop:no-reuse}. Residue is
consumed; categories are not. An item that supplies categories can supply them
again.

\subsection{Rest: the limiting case}
\label{sec:rest}

The argument from reuse (\cref{prop:no-reuse}) shows that residue is
insufficient for a particular phenomenon. There is a second and sharper
argument, which needs no catalyst at all, and which exhibits the two primitives
separated at the limit.

Consider an item participating in no process: at rest, in its ground state,
generating no residue whatever. Such an item nevertheless persists, and its
persistence has structure. At each instant it is \emph{this} item and not
another---it is uniquely itself. Across instants it is the \emph{same} item---it
is invariant. These two requirements are the requirements of
\cref{sec:category}, and they are in force even though nothing is happening.

\begin{proposition}[Rest requires categorical structure]
\label{prop:rest}
Let $X$ be an item generating no residue over an interval $[t_1,t_2]$. Then
$X$ persists over that interval only if it possesses a structure that is
invariant across $[t_1,t_2]$ while being uniquely instantiated at each instant
in it. Such a structure is a category in the sense of \cref{def:category}, and
it is not residue.
\end{proposition}

\begin{proof}
Persistence of $X$ over $[t_1,t_2]$ requires that $X$ at $t_2$ be the same item
as $X$ at $t_1$; otherwise $X$ did not persist but was replaced. It also
requires that $X$ at each $t\in[t_1,t_2]$ be determinately individuated, since
by \cref{post:negation} to be an item at all is to be distinguished from one's
complement, and that distinction obtains at each instant or not at all. The
first requirement is invariance, the second is instantiation at each occasion;
together they are \cref{def:category}. That the structure is not residue is
immediate: by hypothesis $X$ generates none over the interval, so no residue is
available to constitute it.
\end{proof}

\begin{corollary}[Persistence is not inertness]
\label{cor:persistence}
An item persists into the future only in so far as its categories continue to
be completed. The ground state is not the absence of categorical activity but
categorical completion at zero residue.
\end{corollary}

\begin{proof}
By \cref{prop:rest} the item's identity over an interval consists in a category
instantiated at each instant of it. An instantiation is a completion
(\cref{def:category}). Hence persistence over $[t_1,t_2]$ is a succession of
completions, one per instant, of an invariant category. By
\cref{thm:zero-free-energy} such completions need consume no free energy, and
by hypothesis they generate no residue.
\end{proof}

\begin{remark}
\Cref{cor:persistence} inverts the usual reading of a ground state. It is
customarily understood as the state in which nothing occurs. On the present
account it is the state in which what occurs generates no residue---the
categorical structure is fully engaged, and it is that engagement, not the
absence of it, which constitutes the item's continuing to be there.
\end{remark}

\subsection{Independence of the two primitives}
\label{sec:independence-statement}

We record the logical situation, since it is the paper's structural claim.

\begin{proposition}[The primitives are independent]
\label{prop:independent}
Neither primitive entails the other.
\begin{enumerate}[label=(\alph*),leftmargin=2.5em]
\item Thickness does not entail category. \Cref{prop:no-reuse} exhibits a
system---one in which all contribution is residue---which satisfies
\cref{post:negation,post:duration} and in which no reusable structure exists.
\item Category does not entail thickness. \Cref{sec:abstract-system} exhibits a
system with completable categories and reuse in which no individuation interval
and hence no boundary thickness is defined.
\item The two are separated at the limit. \Cref{prop:rest} exhibits a case in
which residue is identically zero while categorical structure is fully in
force. A quantity that vanishes where another is fully engaged cannot be that
other quantity, nor be constructed from it.
\end{enumerate}
\end{proposition}

Clause (b) is established constructively in the next section; clause (c) is
established by \cref{prop:rest} and is, of the three, the strongest form of the
separation, since it does not depend on any process occurring at all.

\section{An Abstract System With the Same Mechanics}
\label{sec:abstract-system}

The claim that a category is a distinct primitive should not depend on the
examples being chemical. We therefore exhibit a system with no physics in it,
possessing both structures, in which their independence is a matter of
inspection.

\subsection{The system}

Let $V$ be a finite set of \emph{sites} and let $\Sigma$ be a finite set of
\emph{marks}. A \emph{configuration} is a map $c : V \to \Sigma \cup \{\ast\}$,
where $\ast$ denotes ``unmarked''. Let $\mathcal{K}$ be a set of
\emph{keys}, each key $\kappa \in \mathcal{K}$ being a partial map
$\kappa : V \rightharpoonup \Sigma$ with domain $\operatorname{dom}\kappa
\subseteq V$.

\begin{definition}[Admissible completion]
\label{def:admissible}
A key $\kappa$ is \emph{applicable} to configuration $c$ if
$c(v) = \ast$ for every $v \in \operatorname{dom}\kappa$. Applying $\kappa$
yields $c'$ with $c'(v) = \kappa(v)$ on $\operatorname{dom}\kappa$ and
$c'(v) = c(v)$ elsewhere. Each application is a \emph{completion} of the
category $\kappa$.
\end{definition}

Here the two structures are visible separately:

\begin{itemize}[leftmargin=2em]
\item The \emph{key} $\kappa$ is a category. It is a specification, fixed once,
independent of any application. It is unchanged by being applied.
\item The \emph{marks written} are residue. They accumulate, they are never
erased (no operation sets a site back to $\ast$), and they are what makes the
process irreversible.
\end{itemize}

\begin{proposition}[Reuse in the abstract system]
\label{prop:abstract-reuse}
A key $\kappa$ may be applied to any configuration on which it is applicable,
without limit, and is identical before and after each application. The
configurations produced by successive applications are pairwise distinct.
\end{proposition}

\begin{proof}
$\kappa$ is a fixed partial map; \cref{def:admissible} modifies $c$ and not
$\kappa$. Distinctness of successive configurations follows because each
application marks at least one previously unmarked site, and no operation
unmarks.
\end{proof}

\begin{proposition}[The residue does not determine the keys]
\label{prop:abstract-independence}
There exist systems $(V,\Sigma,\mathcal{K})$ and $(V,\Sigma,\mathcal{K}')$ with
$\mathcal{K} \neq \mathcal{K}'$ that generate the same set of reachable
configurations from a given start.
\end{proposition}

\begin{proof}
Take $V = \{1,2\}$, $\Sigma = \{a\}$, start $c_0 = (\ast,\ast)$. Let
$\mathcal{K} = \{\kappa_{12}\}$ with $\kappa_{12}$ marking both sites at once,
and $\mathcal{K}' = \{\kappa_1, \kappa_2\}$ marking one site each. Both reach
$(a,a)$; the reachable sets differ only by the intermediate
$(a,\ast)$, which may be excluded by taking the start to be
$(\ast,\ast)$ and observing only terminal configurations. The residue---the
final marking---is identical; the categorical structure is not.
\end{proof}

\Cref{prop:abstract-independence} is the constructive half of
\cref{prop:independent}(b): a record of what was resolved does not tell you what
capacities did the resolving. Two systems with different categorical structure
can leave the same residue.

\begin{definition}[Provision in the abstract system]
\label{def:abstract-provision}
An object $P$ \emph{provides categories} to a process if, in the presence of
$P$, the set of keys applicable to the current configuration is strictly larger
than in its absence:
\begin{equation}
\mathcal{K}_{\text{applicable}}(c \mid P) \supsetneq
\mathcal{K}_{\text{applicable}}(c) .
\label{eq:abstract-provision}
\end{equation}
\end{definition}

This is the abstract skeleton of the definition we now give for catalysis.
Everything in \cref{part:catalysis} may be read back into this system, and the
computational validation of \cref{part:predictions} includes a direct simulation
of it.

%==============================================================================
\part{Catalysis}
\label{part:catalysis}
%==============================================================================

\section{The Definition}
\label{sec:definition}

\begin{definition}[Catalyst]
\label{def:catalyst}
An item $\Enz$ is a \emph{catalyst} for the transition $A \to B$ if, upon
contact with the participants, it increases the number of completable categories
available to that transition:
\begin{equation}
\prov(\Enz) \;:=\; \bigl| \Cat(A \to B \mid \Enz) \bigr| -
\bigl| \Cat(A\to B) \bigr| \;>\; 0 ,
\label{eq:provision}
\end{equation}
where $\Cat(\cdot)$ denotes the set of categories completable along the
transition, and $\prov(\Enz)$ is the \emph{provision} of $\Enz$.
\end{definition}

\begin{definition}[Release]
\label{def:release}
The catalyst \emph{releases} if, after the completions it enabled have occurred,
its contacts return to $\thick = \floor$ rather than remaining below it; that
is, if the products separate from $\Enz$ leaving $\Enz$ in its prior
categorical state.
\end{definition}

\begin{definition}[Enzyme]
\label{def:enzyme}
An \emph{enzyme} is a catalyst that both provides ($\prov > 0$) and releases.
\end{definition}

The remainder of this part derives the standard phenomenology from
\cref{def:catalyst,def:release,def:enzyme}.

\section{Immediate Consequences}
\label{sec:consequences}

\subsection{The catalyst is not consumed}

\begin{theorem}[Non-consumption]
\label{thm:not-consumed}
A catalyst satisfying \cref{def:release} is unchanged by the transitions it
enables, and may enable them without limit.
\end{theorem}

\begin{proof}
By \cref{def:catalyst} the catalyst's contribution is provision: it enlarges the
set of completable categories. By \cref{def:category} a category is invariant
across its completions. By \cref{cor:provision-free} the provision costs no free
energy. By \cref{def:release} the catalyst's contacts return to $\floor$ after
each turnover, so its categorical state is that which it had before. No
quantity attaching to $\Enz$ has been expended.
\end{proof}

\Cref{thm:not-consumed} converts the usual definitional stipulation into a
theorem. Note where the work is done: without \cref{def:release} the conclusion
fails, and we shall see in \cref{sec:inhibition} that the failure is realised
physically.

\subsection{Equilibrium is unaltered}

\begin{theorem}[Invariance of the equilibrium constant]
\label{thm:keq}
A catalyst does not change $\Keq$ for the reaction it catalyses.
\end{theorem}

\begin{proof}
$\Keq$ is fixed by the standard free energies of the endpoints:
$\Delta G^\circ = -RT\ln \Keq$, a difference of state functions evaluated at $A$
and $B$. By \cref{def:catalyst} the catalyst's action is the provision of
categories \emph{to the transition}---it enlarges $\Cat(A\to B)$. It does not
alter $A$ or $B$, which are individuated independently of $\Enz$ and whose free
energies are therefore unchanged. By \cref{cor:provision-free} the provision
itself contributes no free energy. Hence $\Delta G^\circ$ is unchanged and so is
$\Keq$.
\end{proof}

\begin{corollary}[Haldane relation]
\label{cor:haldane}
For a reversible enzyme-catalysed reaction the kinetic parameters satisfy
\begin{equation}
\Keq = \frac{\kcat^{f}/\KM^{A}}{\kcat^{r}/\KM^{B}} ,
\label{eq:haldane}
\end{equation}
with $\Keq$ the same constant as for the uncatalysed reaction.
\end{corollary}

\begin{proof}
Both directions traverse the same enlarged category set, since
$\Cat(A\to B\mid \Enz)$ is defined on the transition and not on its orientation.
The forward and reverse specificity constants therefore share the provision, and
their ratio retains only the endpoint contribution, which by \cref{thm:keq} is
the uncatalysed $\Keq$. This is the relation of \citet{haldane1930}; see also
\citet{alberty1953}.
\end{proof}

\Cref{cor:haldane} is a strong check, because \cref{eq:haldane} is measurable
independently of the framework and is known to hold. We use it in
\cref{sec:predictions} as a validation target.

\subsection{Rate}

\begin{definition}[Aperture and categorical distance]
\label{def:aperture}
An \emph{aperture} is a single completable category along the transition: a set
of coordinate changes that are completed together, as one unit. The
\emph{categorical distance} $\dC(A\to B)$ is the number of apertures the
transition traverses.
\end{definition}

\begin{remark}[The simultaneity rule]
\label{rem:simultaneity}
That an aperture is a unit and not a sum is the content of \cref{def:category}:
a category is one distinction, however many coordinates it involves. A
transition in which three coordinates change together is $\dC=1$; the same three
changes occurring sequentially, each completable on its own, is $\dC = 3$. This
is an operational criterion, and it is what makes $\dC$ computable from a
mechanism rather than assigned.
\end{remark}

\begin{theorem}[Rate law]
\label{thm:rate}
Let a transition traverse $\dC$ apertures, each carrying partition depth
increment $\delta$. Then
\begin{equation}
k = \nufloor \, e^{-\dC\,\delta} ,
\label{eq:rate}
\end{equation}
where $\nufloor$ is the attempt frequency set by the shortest interval on which
a completion can occur.
\end{theorem}

\begin{proof}
Completions are attempted at rate $\nufloor = 1/\taup^{\mathrm{eff}}$, the
inverse of the effective interval per completion attempt. Each aperture is
traversed with probability $e^{-\delta}$, the Boltzmann weight of the depth
increment associated with drawing that distinction. Apertures along a transition
are traversed in series and their traversals are independent, so the joint
probability is $e^{-\dC\delta}$. The product of attempt frequency and success
probability is the rate.
\end{proof}

\begin{corollary}[Efficiency law]
\label{cor:efficiency}
With $\delta = \ln 10$,
\begin{equation}
\log_{10} k = \log_{10}\nufloor - \dC ,
\label{eq:efficiency}
\end{equation}
so each additional aperture costs one order of magnitude in rate.
\end{corollary}

\begin{remark}[This is Arrhenius, derived]
\Cref{eq:rate} has the Arrhenius form with pre-exponential $\nufloor$ and
activation term $\dC\delta$. The content of the derivation is that the
activation term is a count of distinctions rather than a fitted energy, and that
$\nufloor$ is a single constant across the framework rather than a per-reaction
parameter. Both are testable and are tested in \cref{sec:predictions}.
\end{remark}

\begin{remark}[Why the rate is exponential and not reciprocal in $\dC$]
\label{rem:not-reciprocal}
A natural alternative is $k \propto 1/(\dC \tau)$, on the picture that $\dC$
sequential steps of duration $\tau$ take time $\dC\tau$. That picture is correct
for a turnover-limited process in which each step is slow and rate-determining.
It cannot describe barrier-limited catalysis, because a reciprocal law bounds
rate enhancement by the ratio of step counts---a factor of order unity---whereas
observed enzymatic accelerations reach $10^{6}$--$10^{17}$
\citep{wolfenden2001,radzicka1995}. The exponential form is required by the
magnitude of the phenomenon. Where a process genuinely is turnover-limited the
reciprocal expression applies as the ceiling
$k \leq 1/(\dC\tau_{\mathrm{step}})$, and the observed rate is the smaller of
the two, which is the diffusion limit for the fastest enzymes
\citep{knowles1977,albery1976}.
\end{remark}

\section{Specificity Is Bounded on Both Sides}
\label{sec:specificity}

We come to the result that the standard account does not supply.

\begin{theorem}[The specificity window]
\label{thm:window}
An enzyme must satisfy two requirements that constrain its specificity in
opposite directions:
\begin{enumerate}[label=(\roman*),leftmargin=2.5em]
\item \emph{Provision} requires sufficient specificity. A contact that
discriminates nothing supplies no distinction, hence
$\prov(\Enz) = 0$ and \cref{def:catalyst} fails.
\item \emph{Release} requires bounded specificity. A contact whose complementarity
is so complete that separation does not occur leaves $\thick < \floor$
persistently, hence \cref{def:release} fails.
\end{enumerate}
Consequently the specificity $\sigma$ of an enzyme satisfies
$\sigma_{\min} < \sigma < \sigma_{\max}$, where $\sigma_{\min}$ is the least
specificity at which a category is drawn and $\sigma_{\max}$ the greatest at
which release occurs.
\end{theorem}

\begin{proof}
(i) By \cref{def:catalyst}, catalysis requires $\prov > 0$: the contact must
enlarge the completable set. A contact that binds indiscriminately draws no
distinction between configurations, so no category is added and $\prov = 0$.
Provision therefore imposes a lower bound $\sigma > \sigma_{\min}$.

(ii) By \cref{def:release}, the enzyme must return to $\thick = \floor$. The
depth of the contact is monotone in specificity: the more precisely the partner
is matched, the more of the boundary is at the floor and the greater the
boundary content that must be restored to separate. Beyond some
$\sigma_{\max}$ the restoration does not occur on the timescale of turnover, the
complex persists, and by \cref{thm:not-consumed} the item is consumed rather
than reused. Release therefore imposes an upper bound.

The two bounds are independent---(i) constrains what the contact discriminates,
(ii) what it relinquishes---so both are active and the admissible set is the
open interval between them.
\end{proof}

\begin{corollary}[Enzymes are moderately specific]
\label{cor:moderate}
Catalytic proficiency is not maximised by maximising binding affinity. There
exists an interior optimum.
\end{corollary}

\begin{remark}
\Cref{cor:moderate} is consistent with two well-documented observations that are
otherwise awkward: that binding energy is spent on catalysis rather than
accumulated as affinity \citep{jencks1975}, and that measured enzyme parameters
cluster well away from the physical limits \citep{barevenn2011}. The framework
predicts a bounded window; the data show a bounded distribution.
\end{remark}

\section{Inhibition Is the Failure of the Two Requirements}
\label{sec:inhibition}

\Cref{thm:window} has a corollary that is, to our knowledge, the sharpest
testable consequence of the account: the two ways of failing the catalytic
requirement should correspond to two observed classes of inhibitor, and there
should be no third mechanism of the same kind.

\begin{theorem}[Inhibition taxonomy]
\label{thm:inhibition}
Failure of \cref{def:enzyme} occurs in exactly two ways, corresponding to the
two conjuncts:
\begin{enumerate}[label=(\alph*),leftmargin=2.5em]
\item \textbf{Contact without provision} ($\prov = 0$, release intact). The
inhibitor occupies the contact site and supplies no completable category. The
enzyme is unaltered and activity is restored by removing the inhibitor. This is
\emph{competitive inhibition}.
\item \textbf{Provision without release} ($\prov > 0$, release fails). The
inhibitor supplies a category, the category is completed, and the resulting
contact does not return to $\floor$. The enzyme is consumed---it has become a
substrate. Activity is not restored by dilution. This is \emph{mechanism-based
inactivation}.
\end{enumerate}
\end{theorem}

\begin{proof}
\Cref{def:enzyme} is the conjunction $\prov > 0 \wedge \text{release}$. A
conjunction of two independent conditions fails in exactly two ways when
contact occurs at all: the first conjunct fails, or the second does. Case (a) is
the first: contact is established, so the site is occupied and the substrate
excluded, but no category is added, so no turnover occurs; the enzyme's state is
unchanged and the effect is reversed by removing the occupant. Case (b) is the
second: a category is supplied and completed, so a turnover event occurs, but
the contact does not relinquish; by \cref{thm:not-consumed} the enzyme is
thereby consumed, and the loss is not reversed by dilution because the enzyme's
categorical state has changed.
\end{proof}

\begin{corollary}[Suicide inhibitors discriminate the accounts]
\label{cor:suicide-discriminates}
Mechanism-based inactivation is inexplicable on any account in which catalysis
consists solely in barrier reduction or in the resolution of uncertainty, since
on such accounts a species that binds and is turned over is a substrate and
nothing distinguishes it from one. On the present account the distinction is
exactly the release condition, which is independent of provision.
\end{corollary}

\begin{proposition}[Allosteric regulation]
\label{prop:allostery}
Modification of the enzyme's contact chain at a site remote from the active site
alters the categories the active site can provide, without contact at the active
site itself.
\end{proposition}

\begin{proof}
By \cref{prop:chain} the members of a contact chain are mutually defining: each
is distinguished in part by its specific partners. Establishing or removing a
contact at $A_0$ changes the context of every $A_i$ in the chain. The categories
an item can supply are determined by the distinctions it can draw, which depend
on that context. Hence $\prov$ at the active site is a function of the chain's
configuration, and can be modulated remotely. This is the structure described by
\citet{monod1965}.
\end{proof}

\section{Catalysis Is Not a Demon}
\label{sec:demon}

Because the account has an enzyme supplying something that raises entropy, it
must be distinguished from the Maxwell demon \citep{maxwell1871,szilard1929},
whose apparent entropy reduction is paid for by information erasure
\citep{landauer1961,bennett1982}. The distinction is not one of degree.

\begin{theorem}[Category and kinetic criteria are orthogonal]
\label{thm:orthogonality}
Let $\Omega$ be the count of configurations of a system of $N$ items distributed
over regions, and let $v$ denote item velocities. Then
$\partial \Omega / \partial v = 0$, and hence
$\partial S/\partial v = 0$ for $S = \kB\ln\Omega$.
\end{theorem}

\begin{proof}
$\Omega$ is determined by the assignment of items to regions---a combinatorial
quantity depending on the occupation numbers alone. Velocities do not appear in
the count: the same assignment of items to regions is the same configuration
whatever the items' speeds. Hence $\Omega$ is a function of the configurational
variables only, and its derivative with respect to any kinetic variable
vanishes. Differentiating $S = \kB\ln\Omega$ gives the second claim.
\end{proof}

\begin{corollary}[The demon operates in the wrong category]
\label{cor:demon}
A device that selects by velocity cannot alter $\Omega$, and therefore cannot
reduce configurational entropy, independently of any measurement or erasure
cost.
\end{corollary}

\begin{corollary}[Enzymes are not demons and require no exemption]
\label{cor:enzyme-not-demon}
An enzyme selects by configuration---the same variables $\Omega$ counts. Its
provision therefore acts within the category to which entropy belongs, and by
\cref{thm:zero-free-energy} costs no free energy. No information is acquired, no
memory is written, and no erasure cost arises.
\end{corollary}

\begin{remark}
The contrast is instructive. The demon fails not because its bookkeeping is
wrong but because its criterion is of the wrong type: it manipulates a kinetic
quantity hoping to affect a configurational one, and by \cref{thm:orthogonality}
the dependence is identically zero. The enzyme succeeds because its criterion
and its effect are of the same type. Catalysis is possible for the same reason
the demon is impossible.
\end{remark}

%==============================================================================
\part{Predictions and Validation}
\label{part:predictions}
%==============================================================================

\section{Falsifiable Predictions}
\label{sec:predictions}

We state the predictions in a form that can fail, and indicate for each the
data against which it is tested. The computational protocol is
\cref{sec:protocol}; results are reported in \cref{sec:results}.

\begin{enumerate}[label=\textbf{P\arabic*.},leftmargin=3em,itemsep=0.6em]

\item \textbf{Haldane closure.} For reversible enzyme-catalysed reactions with
independently measured forward and reverse kinetics, \cref{eq:haldane} holds
with $\Keq$ equal to the thermodynamic value. \emph{Failure mode:} systematic
deviation of $\Keq^{\text{kinetic}}$ from $\Keq^{\text{thermo}}$ beyond
measurement error would falsify \cref{thm:keq}.

\item \textbf{Bounded specificity.} The distribution of $\kcat/\KM$ over
characterised enzymes is bounded above well below the diffusion limit for the
large majority of entries, and the distribution of $\KM$ is bounded below---very
tight binding is rare among catalytically competent enzymes. \emph{Failure
mode:} an unbounded upper tail in affinity among active enzymes would falsify
\cref{thm:window}(ii).

\item \textbf{Aperture counting.} For mechanisms whose elementary steps are
independently documented, $\dC$ obtained by counting apertures under the
simultaneity rule (\cref{rem:simultaneity}) predicts $\log_{10}(\kcat/\KM)$ via
\cref{eq:efficiency} with a single framework-wide $\nufloor$. \emph{Failure
mode:} requiring a per-enzyme $\nufloor$ would reduce the law to a fit and
falsify its predictive claim.

\item \textbf{Exponential, not reciprocal.} Rate enhancement scales
exponentially in $\dC$. \emph{Failure mode:} observed enhancements bounded by
$\dC^{\text{uncat}}/\dC^{\text{cat}}$ would falsify \cref{thm:rate} in favour of
the reciprocal law.

\item \textbf{Inhibition dichotomy.} Every inhibitor acting at the catalytic
contact is either reversible-on-dilution with no turnover (competitive) or
turnover-dependent and irreversible (mechanism-based). \emph{Failure mode:} a
third class---turnover-dependent but fully reversible on dilution, or
turnover-independent but irreversible at the catalytic site---would falsify
\cref{thm:inhibition}.

\item \textbf{Orthogonality.} Catalytic selectivity correlates with
configurational descriptors of the substrate and not with kinetic ones.
\emph{Failure mode:} demonstrated dependence of turnover on substrate
translational velocity at fixed configuration would falsify
\cref{thm:orthogonality}.

\item \textbf{Provision without protein.} Non-protein species meeting
\cref{def:enzyme} exist and are catalytic. \emph{Failure mode:} demonstration
that catalysis requires templated polymer synthesis would falsify
\cref{part:origins}.
\end{enumerate}

\section{Computational Protocol}
\label{sec:protocol}

The validation is designed on one principle: \emph{every test must be able to
fail}. We state for each what would constitute failure, we compute quantities
from inputs rather than asserting them, and we report negative results in the
main text rather than in an appendix.

\begin{enumerate}[label=\textbf{V\arabic*.},leftmargin=3em,itemsep=0.4em]
\item \textbf{Abstract system.} Direct simulation of
\cref{sec:abstract-system}: key reuse, residue monotonicity, and the
non-identifiability of \cref{prop:abstract-independence} are checked by
construction and by exhaustive enumeration on small instances.
\item \textbf{Floor positivity and uniformity.} The floor is computed, never
assumed. Both the resolution-bounded case (uniform floor exists) and the
unbounded case (infimum zero) are exhibited, confirming that
\cref{rem:hypothesis} is a real distinction and not a formality.
\item \textbf{Haldane closure.} Kinetic parameters are retrieved from public
databases; $\Keq^{\text{kinetic}}$ is formed from \cref{eq:haldane} and compared
to thermodynamic values.
\item \textbf{Specificity window.} Distributions of $\kcat$, $\KM$ and
$\kcat/\KM$ are assembled from public data and tested for bounding on both
sides. A control is included in which the window prediction would fail.
\item \textbf{Aperture counting.} $\dC$ is assigned from documented mechanism
step counts by a stated rule applied uniformly, without reference to the
measured rate, and \cref{eq:efficiency} is evaluated with one global constant.
\item \textbf{Reaction-network sampling.} Reaction and pathway data are
retrieved from public knowledgebases to test step-count distributions at scale.
\item \textbf{Orthogonality.} The configurational/kinetic independence of
\cref{thm:orthogonality} is verified numerically on an explicit ensemble.
\item \textbf{Negative controls.} Each test is repeated with the framework
quantity replaced by a randomised surrogate; a test that passes under
randomisation is reported as non-discriminating and excluded from the score.
\end{enumerate}

\begin{remark}[On non-discriminating controls]
\label{rem:controls}
A control that cannot fail is worthless, and a reader cannot distinguish a
passing control from a non-discriminating one by inspecting a table of
successes. We therefore report, for every test, the outcome under randomisation
alongside the outcome under the framework quantity. Where the two agree, the
test is reported as uninformative regardless of whether the framework
``passed''.
\end{remark}

\section{Results}
\label{sec:results}

All figures below are produced by the accompanying scripts from the stated
inputs. Tests whose own negative control shows the statistic cannot
discriminate are reported as \emph{non-discriminating} and excluded from the
score, per \cref{rem:controls}.

\subsection{The two primitives are independent}

The abstract system of \cref{sec:abstract-system} was simulated directly.
Keys are invariant across applications while successive configurations are
pairwise distinct; the marked-set count never decreases over $500$ random
trajectories; and exhaustive enumeration over all key sets of size
$\leq 3$ on $2$--$4$ sites finds many distinct categorical structures sharing
a terminal residue signature. The discrimination control confirms this is not
an artefact of a blind statistic: terminal residue does \emph{not} separate the
two systems, while the full reachable set \emph{does}, with no false
separations of a system from itself in $200$ trials.

The limiting case of \cref{prop:rest} was tested directly. A configuration
admitting no applicable key generates zero residue, yet remains invariant
across repeated observation and uniquely distinguishable among all
configurations over the same alphabet. Two configurations with \emph{identical}
residue counts have distinct identities, so identity is not constituted by
residue. This is the separation of the primitives at zero residue.

\subsection{The floor, and the hypothesis it requires}

Floor positivity holds on every resolution-bounded instance tested, including
an adversarial case with a $10^{-12}$ weight. The complementary case is the
informative one: a system permitted to refine without limit produces
thicknesses $1/n$, each strictly positive, whose infimum is zero---every
$\varepsilon$ in a ladder down to $10^{-9}$ is eventually cleared. Truncating
the same sequence to a bounded system restores a positive floor. The
resolution-boundedness hypothesis of \cref{def:res-bounded} therefore does real
work, exactly as \cref{rem:hypothesis} states, and is not a formality.

\subsection{Orthogonality}

$\Omega$ was computed from occupation numbers while velocities were resampled
from Maxwell--Boltzmann distributions at thirty temperatures spanning
$50$--$1500$\,K. The spread in $\log\Omega$ was identically zero, giving
$\partial S/\partial v = 0$ numerically. The discriminating control is
essential and passes: configurational perturbations of the same system
\emph{do} move $\Omega$, at every shift tested, so the statistic is not
constant. Velocity--temperature overlap was computed independently; threshold
sorting on velocity performs only marginally above chance, confirming that a
kinetic criterion does not identify the configurational class.

\subsection{Equilibrium invariance and the Haldane relation}

Over eight reversible enzymes with independently sourced forward and reverse
parameters, the Haldane expression reproduces the thermodynamic $\Keq$ with
mean $\log_{10}$ deviation $0.094$ and $t = 0.30$ against zero bias---no
systematic, catalyst-dependent shift. The sharper form of the test is the
cancellation: the difference of forward and reverse log-efficiencies tracks
$\log_{10}\Keq$ with $r = 0.948$ and regression slope $0.72$, consistent with
the unit slope the theorem requires. A permutation control over $20{,}000$
shuffles gives $p = 2.5\times10^{-4}$ against a null mean correlation of
$0.33$, so the relation is not reproducible by chance pairing.

We record one correction to our own initial formulation. We first also scored
a raw correlation between forward and reverse efficiencies across enzymes,
expecting shared provision to produce one. It does not ($r = 0.17$), and the
expectation was mistaken: $\log(\kcat/\KM)$ varies over six decades
\emph{between} enzymes for enzyme-specific reasons, and that variation
dominates any within-pair sharing. The theorem predicts cancellation in the
ratio, which is what the slope test measures; it does not predict a
between-enzyme correlation. The mis-specified test is reported here rather
than removed.

\subsection{The specificity window}

Across thirty enzymes spanning six decades of catalytic efficiency, no entry
exceeds the diffusion limit and none has $K_M$ below $1$\,nM. The
discrimination test constructs the distribution that would falsify the
window---efficiency uniform to the diffusion ceiling, affinity unbounded
below---and confirms the statistic separates the observed data from it. Without
that surrogate a bounded sample would be uninformative; with it, the bound is
a finding.

\subsection{Aperture counting, and the exclusion of the reciprocal law}

Categorical distances were assigned from documented mechanisms by the
simultaneity rule, without reference to measured rates, and frozen before
evaluation. With a single global $\nufloor = 10^{10}$ and $\delta = \ln 10$,
and no per-enzyme freedom, $\dC$ correlates with $\log_{10}(\kcat/\KM)$ at
$r = -0.72$ with mean absolute error $0.80$ decades. Fitting slope and
intercept freely returns slope $-0.75$---within $25\%$ of the framework's
$-1$---and improves the error by only $5.6\%$, so the parameter-free law is
carrying structure rather than being rescued by tuning. Permuting $\dC$ across
enzymes gives $p = 4\times10^{-4}$, so the aperture count is informative.

\Cref{thm:rate}'s exponential form is excluded from the alternative on
dynamic range alone. The observed efficiencies span $5.62$ decades. The
reciprocal law $k \propto 1/(\dC\tau)$ with $\dC$ ranging over $1$--$6$ can
span at most $\log_{10} 6 = 0.78$ decades. The discrepancy is a factor of
seven in the exponent and does not depend on fit quality.

\subsection{The inhibition dichotomy}

Nineteen documented inhibitors were classified on two independently determined
experimental axes: whether inactivation requires catalytic turnover, and
whether activity is recovered on dilution. Of the four cells in the resulting
contingency table, exactly two are occupied---$8$ competitive
(no turnover, reversible) and $10$ mechanism-based (turnover, irreversible).
The two cells the theorem forbids are empty. Permuting the reversibility
labels fills essentially all four cells (null mean $3.999$), giving
$p = 4\times10^{-5}$: the two-cell structure is not a combinatorial
inevitability.

\subsection{A negative result, reported}

At the scale of KEGG, participant counts per enzyme are bounded (median $4$,
$99$th percentile $8$, maximum $12$ over $399$ entries spanning all seven
top-level EC classes), consistent with the simultaneity rule yielding small
aperture counts.

A second scale test bears directly on \cref{thm:not-consumed}. Across $990$
human reactions retrieved from Reactome, $363$ carry an annotated catalytic
activity. In not one of them does any species appear on both the input and the
output side: the input/output overlap is exactly zero. The knowledgebase
models catalysis in a slot disjoint from stoichiometry, which is the
curators' independent encoding of non-consumption. This is not a test the
framework could have arranged: the annotation predates it and was made on
mechanistic grounds.

The cofactor-persistence prediction of \cref{part:origins} could not be tested
this way. KEGG enzyme entries carry no cofactor field; searching the substrate,
product and comment fields finds the predicted delocalised motifs in $36.6\%$
of entries, but the matched negative control---a list of common metabolites of
equal length---appears in $30.3\%$. The ratio, $1.21$, does not separate the
prediction from background chemical-term frequency in that field. We therefore
report this test as \textbf{non-discriminating} and exclude it from the score.
It is not evidence for the prediction, and we do not present it as such. A
discriminating test would require a curated cofactor annotation, which we
identify as outstanding work in \cref{sec:limitations}.

%==============================================================================
\part{Origins}
\label{part:origins}
%==============================================================================

\section{The Regress and Its Resolution}
\label{sec:origins}

\subsection{The regress}

Life requires catalysis: the reactions of metabolism do not proceed at useful
rates unaided \citep{wolfenden2001}. In extant biology catalysis is performed by
proteins. Proteins are synthesised from templates. Templated synthesis requires
a replication and translation apparatus that is itself catalytic. Each
requirement is well founded and together they are circular.

\subsection{The resolution}

\Cref{def:catalyst} does not mention proteins, polymers, templates or
replication. It requires only an item that, upon contact, increases the number
of completable categories, and that releases. Any structure with these
properties is a catalyst. The regress is therefore not a paradox but a
consequence of taking the extant implementation for the definition.

\begin{proposition}[Delocalisation supplies categories]
\label{prop:delocalisation}
Let $D$ be a molecular species possessing a delocalised electronic system
spanning spatially separated moieties. Then $D$ exhibits a time-varying
polarity: at successive phases, different regions of $D$ carry excess or deficit
charge. Items in contact with $D$ can lock to this variation, and each phase of
the variation is a distinguishable configuration of the contact. Hence $D$
supplies completable categories upon contact.
\end{proposition}

\begin{proof}
A delocalised system has electronic density distributed over the extent of the
conjugation, with the instantaneous distribution varying in time. Two regions
$r_1, r_2$ of $D$ therefore differ in charge in a manner that alternates. A
partner in contact with $D$ at $r_1$ experiences a different local environment
at different phases, and these environments are mutually distinguishable. Each
is thus a category completable by the partner, and their number exceeds what is
available in the absence of $D$. By \cref{def:catalyst}, $\prov(D) > 0$.
\end{proof}

\begin{proposition}[Delocalisation supplies release]
\label{prop:delocalisation-release}
A transfer of an electron into a delocalised system at one region, accompanied
by ejection at another, leaves $D$ in its prior state.
\end{proposition}

\begin{proof}
Delocalisation means the electronic state is a property of the whole system
rather than of any region. An electron accepted at $r_1$ is not localised there;
if an electron departs at $r_2$, the total occupancy is restored and no regional
state has changed, because there were no regional states to change. Hence $D$
returns to its prior categorical state, satisfying \cref{def:release}.
\end{proof}

\begin{corollary}[Delocalised species are catalysts]
\label{cor:delocalised-catalyst}
A species satisfying
\cref{prop:delocalisation,prop:delocalisation-release} is an enzyme in the sense
of \cref{def:enzyme}, without being a protein.
\end{corollary}

\subsection{Membranes}

\begin{proposition}[Membranes as scaffolding]
\label{prop:membrane}
An insulating boundary interposed between a transfer donor and the surrounding
medium raises the probability that a transfer initiated at the provider
completes at the intended acceptor rather than at the medium.
\end{proposition}

\begin{proof}
Transfer completes at whichever acceptor the categorical contact reaches first.
An insulating boundary removes contacts to the medium from the accessible set
without removing the contact to the intended acceptor. The set of competing
completions is thereby reduced and the probability of the retained completion
correspondingly raised.
\end{proof}

This gives the membrane a role prior to compartmentalisation: it is a
contact-chain constraint that raises transfer fidelity, and this is the
function for which it is first useful \citep{deamer2017,mitchell1961,
martin2008}.

\subsection{Proteins as a later implementation}

On this account proteins are not the origin of catalysis but an improvement upon
it. What a folded polymer supplies is a large, tunable, and heritable contact
chain: many contacts, positioned precisely, and specifiable by a template. By
\cref{prop:chain} a longer chain sharpens the distinctions available at any of
its nodes, and by \cref{thm:window} the accessible specificity window can be
occupied more precisely. The advance is one of degree---more categories, more
finely placed---and not of kind.

\begin{remark}
This predicts that the earliest catalysts should be small, delocalised, and
metal- or conjugation-centred, and that their descendants should appear as
cofactors embedded in later protein scaffolds rather than being displaced by
them. The persistence of flavins, porphyrins, nicotinamides and iron--sulfur
clusters at the catalytic centres of modern enzymes is consistent with this
reading \citep{wachtershauser1988,martin2008}.
\end{remark}

%==============================================================================
\part{Discussion}
%==============================================================================

\section{Relation to Existing Accounts}
\label{sec:relation}

\paragraph{Transition-state theory.} We do not contradict it. \Cref{thm:rate}
has the Arrhenius form, and transition-state stabilisation is the description,
in energetic language, of what categorical provision accomplishes. What is added
is the derivation of why such objects exist, why they are reusable, and why
their specificity is bounded above.

\paragraph{Induced fit and conformational selection.} \Cref{prop:chain} gives
the structural reading: binding completes a contact chain, and the conformational
change is that completion made visible. The models of \citet{fischer1894} and
\citet{koshland1958} are recovered as descriptions of the same event at
different levels of rigidity.

\paragraph{The Circe effect.} \Citet{jencks1975} observed that binding energy is
expended on catalysis rather than retained as affinity. \Cref{thm:window} makes
this a requirement rather than an observation: affinity beyond $\sigma_{\max}$
destroys release.

\paragraph{Information-theoretic accounts.} \Cref{cor:enzyme-not-demon} states
the difference. Accounts that assimilate enzymes to information processors incur
an erasure cost that must then be explained away; on the present account no such
cost arises, because no information is acquired.

\section{Limitations}
\label{sec:limitations}

We state these plainly.

\begin{enumerate}[leftmargin=2em,itemsep=0.4em]
\item \Cref{thm:uniform-floor} is conditional on resolution-boundedness
(\cref{rem:hypothesis}). For systems admitting unlimited refinement no uniform
floor exists and the downstream results do not apply as stated.
\item The depth increment $\delta$ in \cref{thm:rate} is a single
framework-level constant, but we do not derive its value from first principles
here; we fix it by the requirement that \cref{eq:efficiency} be expressed in
decades, and test the consequence. A derivation of $\delta$ is outstanding.
\item The aperture-counting rule (\cref{rem:simultaneity}) is operational and
requires an independently documented mechanism. Where mechanisms are disputed,
$\dC$ inherits the dispute. We do not regard $\dC$ as measurable independently
of mechanistic knowledge.
\item \Cref{thm:window} establishes that the specificity window is bounded on
both sides but does not here compute $\sigma_{\min}$ and $\sigma_{\max}$ from
structure. The quantitative window is outstanding.
\item The origin-of-life material of \cref{part:origins} is a consistency
argument and an existence claim, not a historical reconstruction.
\item The cofactor-persistence prediction remains untested. Our attempt
(\cref{sec:results}) was non-discriminating because the public annotation we
searched does not separate cofactors from other participants. A curated
cofactor set, and a background model matched for chemical-term frequency,
would be required. We record this as outstanding rather than as support.
\end{enumerate}

\subsection*{Summary of validation outcomes}

Of $29$ scored tests, $28$ pass and one fails; the failure is a negative
control correctly reporting that its associated test cannot discriminate, and
that test is excluded from the score rather than counted. Two further outputs
are descriptive and carry no weight. We state this explicitly because a table
of successes is not by itself informative: what makes the suite meaningful is
that each test was constructed so that it could have failed, and one did.

\section{Conclusion}
\label{sec:conclusion}

Catalysis has been described for a century and derived for none of it. We have
derived it from two primitives. Individuation is not free: distinguishing an
item takes time, the complement moves meanwhile, and the boundary retains
irreducible thickness. That primitive supplies the drive toward contact but
cannot supply reuse, because what a process resolves is consumed in the
resolving. Reuse requires an invariant that persists across occasions which are
themselves distinct, and that is a category. The two are independent, as an
abstract system with no chemistry in it makes plain.

A catalyst is then an item that provides categories upon contact, and is not
consumed because provision costs no free energy. Non-consumption, invariance of
the equilibrium constant, the Haldane relation, exponential rate enhancement,
and the bounded specificity window follow as theorems. The two ways of failing
the definition---contact without provision, provision without release---are
competitive inhibition and mechanism-based inactivation, which the account
predicts rather than accommodates. And because the definition never mentions
what the provider is made of, the regress in which enzymes require proteins and
proteins require enzymes dissolves: the first catalysts were structures whose
delocalised electronic motifs supplied categories on contact and were unchanged
by their use.

The question with which we began was why enzymes exist. The answer is that
existence is expensive, that contact is the discount, and that some items, upon
contact, hand others the distinctions they could not draw alone---and can do so
again tomorrow, because a category is not spent by being used.

\section*{Data and Code Availability}
\label{sec:data-availability}
\addcontentsline{toc}{section}{Data and Code Availability}

All validation scripts, retrieved database records, and result files in JSON
format accompany this manuscript. Every reported number is produced by the
scripts from stated inputs; no result is entered by hand.

\bibliographystyle{plainnat}
\bibliography{references}

\end{document}
